\documentclass[a4paper,10pt]{article}
\usepackage[utf8]{inputenc}
\usepackage[T1]{fontenc}
\usepackage[margin=2cm]{geometry}
\usepackage{xcolor}
\usepackage{hyperref}
\usepackage{tabularx}
\usepackage{amsmath}
\usepackage{array}

\hypersetup{
    colorlinks=true,
    linkcolor=blue!70!black,
    urlcolor=blue!70!black,
    citecolor=blue!70!black,
    pdftitle={Strudel Studio - Manuel Utilisateur et Guide Complet},
    pdfauthor={Hugo Birard}
}

\definecolor{codegray}{gray}{0.95}
\definecolor{primary}{RGB}{99, 102, 241}
\definecolor{accent}{RGB}{16, 185, 129}
\definecolor{darkbg}{RGB}{15, 23, 42}

\setlength{\parindent}{0pt}
\setlength{\parskip}{6pt}

\begin{document}

\begin{titlepage}
    \centering
    \vspace*{2cm}
    {\Huge \textbf{\color{primary} STRUDEL STUDIO}}\\[0.5cm]
    {\Large \textbf{Application Desktop Dédiée au Live-Coding Musical}}\\[0.3cm]
    {\large Manuel Officiel d'Utilisation \& Guide de Référence Exhaustif}\\[1.5cm]

    \rule{\linewidth}{0.5mm}\\[0.4cm]
    {\large \textbf{Auteur :}} Hugo Birard\\
    {\large \textbf{Moteur Audio :}} Strudel / TidalCycles WebAudio Architecture\\
    {\large \textbf{Version :}} 2.0 (Architecture Modulaire \& CodeMirror 6)\\[0.4cm]
    \rule{\linewidth}{0.5mm}\\[2cm]

    \begin{quote}
    \centering
    \textit{Un environnement de live-coding sans compromis : interface sombre sur-mesure, bus matériel WebAudio réel, éditeur professionnel à coloration syntaxique, métronome zero-drift et bibliothèque personnelle de snippets.}
    \end{quote}

    \vfill
    {\large Novembre 2026}
\end{titlepage}

\tableofcontents
\newpage

\section{Introduction \& Philosophie}
\textbf{Strudel Studio} est une application pour ordinateur de bureau (Linux, macOS, Windows) conçue pour la création musicale algorithmique en temps réel. Elle intègre le moteur de synthèse et de séquençage de \textbf{Strudel} (implémentation JavaScript de \textbf{TidalCycles}), tout en remplaçant l'interface web par défaut par un studio de production épuré, sombre, modulaire et hautement performant.

\subsection{Principes Directeurs}
\begin{enumerate}
    \item \textbf{Zéro code imposé :} L'éditeur s'ouvre vierge, sans aucun exemple pré-rempli. Vous démarrez toujours d'une page blanche.
    \item \textbf{Flux matériel réel :} Les curseurs de réglage (Volume Master, Filtre Passe-Bas Master) sont physiquement connectés au graphe de traitement WebAudio natif.
    \item \textbf{Édition professionnelle :} Intégration officielle de l'éditeur \textbf{CodeMirror 6} avec coloration syntaxique complète du langage Strudel et de la mini-notation Tidal.
    \item \textbf{Synchronisation parfaite :} Asservissement de l'affichage métrique sur l'horloge interne de Strudel pour éliminer toute dérive temporelle.
    \item \textbf{Liberté totale :} Bibliothèque de snippets personnalisée persistante dans le stockage local de l'utilisateur.
\end{enumerate}

\section{Guide d'Installation \& Lancement}

\subsection{Prérequis Système}
\begin{itemize}
    \item \textbf{Système d'exploitation :} Linux (Ubuntu, Debian, Fedora, Arch, etc.), Windows 10/11 ou macOS.
    \item \textbf{Node.js :} Version 18 ou supérieure recommandée (\texttt{node -v}).
    \item \textbf{Gestionnaire de paquets :} npm ou pnpm.
    \item \textbf{Carte son \& Pilotes WebAudio :} Compatible ALSA, PulseAudio ou PipeWire sous Linux.
\end{itemize}

\subsection{Installation depuis le Dépôt Git}
Ouvrez un terminal et exécutez les commandes suivantes :
\begin{verbatim}
# 1. Cloner le dépôt officiel
git clone git@github.com:HugoBirard/Strudel.git
cd Strudel

# 2. Installer les dépendances du projet
npm install

# 3. Lancer l'application
npm start
\end{verbatim}

\subsection{Installation du Raccourci Système (Bureau \& Menu d'applications)}
Pour lancer Strudel Studio en un clic sans ouvrir de terminal, exécutez le script d'installation du raccourci :
\begin{verbatim}
npm run install-shortcut
# ou directement :
bash scripts/install-shortcut.sh
\end{verbatim}
Ce script crée une entrée de bureau sur votre \textbf{Bureau} (\texttt{\textasciitilde/Desktop} ou \texttt{\textasciitilde/Bureau}) ainsi que dans votre lanceur d'applications Linux (\texttt{\textasciitilde/.local/share/applications/strudel-studio.desktop}).

\subsection{Compilation d'un Binaire Autonome (AppImage Linux)}
Pour packager l'application sous la forme d'un exécutable unique transportable sans nécessiter Node.js :
\begin{verbatim}
npm run build:linux
\end{verbatim}
Le fichier généré sera disponible dans le répertoire \texttt{dist/} sous le nom \texttt{Strudel Studio-1.0.0.AppImage}.

\newpage

\section{Visite Guidée de l'Interface Utilisateur}
L'interface de Strudel Studio a été conçue pour éliminer toute friction visuelle, allier esthétique sombre et ergonomie de production, et s'adapter avec fluidité à tous les agencements d'écrans (du plein écran 4K jusqu'au mode scindé côte-à-côte).

\subsection{1. Barre de Titre Système CSD \& Contrôles de Fenêtre (Style ToDoList)}
Strudel Studio intègre une barre de titre personnalisée \textit{Client-Side Decorations} (CSD) sans bordure native encombrante :
\begin{itemize}
    \item \textbf{Contrôles de Fenêtre Intégrés :}
    \begin{itemize}
        \item \textbf{Réduire (\texttt{[-]}) :} Minimise l'application dans la barre des tâches de votre environnement de bureau.
        \item \textbf{Agrandir / Restaurer (\texttt{[+]} / \texttt{[Restaurer]}) :} Alterne instantanément entre l'affichage plein écran et la dimension normale restaurée. Cette fonction est essentielle pour disposer confortablement Strudel Studio à côté d'une autre fenêtre (navigateur web, éditeur tiers, terminal).
        \item \textbf{Fermer (\texttt{[x]}) :} Ferme proprement l'application et libère l'AudioContext.
        \item \textbf{Double-clic rapide :} Un double-clic sur la barre de titre bascule également l'état maximisé / restauré.
    \end{itemize}
    \item \textbf{Redimensionnement Fluide 8 Directions :} La fenêtre intègre 8 zones de détection actives aux 4 bordures et 4 coins (\texttt{n}, \texttt{s}, \texttt{e}, \texttt{w}, \texttt{ne}, \texttt{nw}, \texttt{se}, \texttt{sw}) permettant d'ajuster librement la taille de la fenêtre à la souris sous n'importe quel gestionnaire de fenêtres (GNOME Wayland/X11, KDE Plasma, Windows 10/11, gestionnaires de fenêtres en tuiles).
    \item \textbf{5 Styles Visuels Personnalisables :} Accessibles depuis le sélecteur de thèmes :
    \begin{enumerate}
        \item \texttt{match-bg} \textit{(Par défaut)} : Aspect monolithique uni, parfaitement fusionné avec la nuance sombre du thème sélectionné.
        \item \texttt{glass} : Effet de verre translucide \textit{Glassmorphism} avec flou d'arrière-plan, reflet supérieur et ombre portée.
        \item \texttt{neon} : Liseré supérieur cyberpunk illuminé avec halo indigo phosphorescent.
        \item \texttt{transparent} : Transparence intégrale sans séparation visuelle (style Spotify / Discord).
        \item \texttt{minimal} : Hauteur ultra-réduite (28 px) masquant le titre pour dédier 100\% de l'espace au code.
    \end{enumerate}
\end{itemize}

\subsection{2. Barre de Transport \& Synchronisation Métrique Zero-Drift}
\begin{itemize}
    \item \textbf{Bouton PLAY (\texttt{Ctrl + Entrée}) :} Lance l'évaluation du code en cours dans l'éditeur. Réveille instantanément l'AudioContext en cas de restriction d'autoplay du système d'exploitation.
    \item \textbf{Bouton STOP (\texttt{Ctrl + .}) :} Coupe immédiatement tous les cycles sonores et libère les oscillateurs actifs.
    \item \textbf{Métronome Visuel 4 Temps (Zero-Drift) :} 4 diodes lumineuses indiquent en temps réel les pulsations du cycle rythmique. Contrairement aux métronomes classiques utilisant des timers JavaScript décalés, ce métronome est directement asservi à l'horloge haute précision de Strudel (\texttt{repl.scheduler.now()}).
    \item \textbf{Contrôleur de Tempo (BPM) \& TAP :} Permet d'ajuster le tempo à la volée avec les boutons $+$ / $-$, par saisie numérique ou en battant la mesure sur le bouton \textbf{TAP}.
    \item \textbf{Bouton Thème :} Ouvre un panneau permettant de basculer instantanément entre 6 thèmes d'interface personnalisés et de sélectionner le style de barre de titre CSD.
    \item \textbf{Boutons Fichier (Nouveau, Ouvrir, Sauvegarder) :} Permettent de gérer vos fichiers de patterns musicaux locaux (\texttt{.strudel}, \texttt{.js}).
\end{itemize}

\subsection{3. Panneau Gauche : Bibliothèque de Snippets \& Tiroir Rétractable}
Ce panneau permet de conserver vos briques de composition favorites :
\begin{itemize}
    \item \textbf{Préréglages Immédiats :} Prise en main instantanée avec des patterns prêts à l'emploi (Kick 808/909, Bass Acid 303, Nappe Rhodes, Arpège d'ambiance).
    \item \textbf{Insertion chirurgicale au curseur :} Cliquer sur \textit{"Insérer au curseur"} injecte exactement le code à l'emplacement actif du curseur dans l'éditeur, sans ajouter de saut de ligne parasite.
    \item \textbf{Bouton "Sauvegarder la sélection" :} Sélectionnez une portion de texte dans l'éditeur à la souris et cliquez sur ce bouton pour créer un nouveau snippet pré-rempli.
    \item \textbf{Mode Tiroir en Écran Partagé :} En largeur compacte, le panneau se replie et s'ouvre sous forme de tiroir fluide en surimpression d'un clic sur le bouton d'en-tête \texttt{[Snippets]} sans réduire la surface du code.
\end{itemize}

\subsection{4. Panneau Central : Éditeur Professionnel CodeMirror 6 \& Architecture Split-Screen}
Le cœur de travail de l'application bénéficie de toute la puissance de CodeMirror 6 :
\begin{itemize}
    \item \textbf{Conception Responsive Split-Screen (Largeur Mini 480 px) :}
    L'application accepte un redimensionnement jusqu'à $480 \times 360$ pixels pour permettre une juxtaposition fluide côte-à-côte avec un autre outil de développement, un navigateur web ou la documentation.
    En mode compact ($\le 920$ px), \textbf{le code reste 100\% visible et manipulable} grâce au défilement vertical et horizontal CodeMirror (\texttt{.cm-scroller}) qui bénéficie de toute la hauteur de fenêtre disponible.
    \item \textbf{Démarrage Sonore Immédiat :} Prévention active des restrictions d'autoplay du navigateur et mise en cache des échantillons de percussion (\texttt{RolandTR909}). L'application démarre avec un motif rythmique et une ligne de basse prête à jouer.
    \item \textbf{Coloration syntaxique complète :} Reconnaissance native des fonctions Strudel (\texttt{s}, \texttt{note}, \texttt{stack}, \texttt{cpm}), de la mini-notation TidalCycles (crochets, chevrons, alternances, accélérations), des variables et des nombres.
    \item \textbf{Auto-fermeture \& Indentation :} Fermeture automatique des délimiteurs et gestion intelligente de la touche Tabulation.
    \item \textbf{Barre de statut dynamique :} Affiche en temps réel si le moteur audio est à l'arrêt, en cours d'évaluation, ou si une erreur de syntaxe a été détectée dans le code.
\end{itemize}

\subsection{5. Panneau Droit : Visualisateur Audio 60 FPS \& Bus Matériel WebAudio}
\begin{itemize}
    \item \textbf{Visualisateur Audio Réactif 60 FPS :}
    \begin{itemize}
        \item \textit{Oscilloscope temporel :} Tracé d'onde illuminé réagissant directement aux formes d'ondes générées (dents de scie, sinus, impulsions).
        \item \textit{Spectre fréquentiel FFT :} Barres d'analyse spectrale animées en arrière-plan révélant la répartition des fréquences basses, médiums et aiguës.
    \end{itemize}
    \item \textbf{Curseur Master Cutoff (Filtre LPF) :} Contrôle en temps réel la fréquence de coupure d'un filtre physique \texttt{BiquadFilterNode} placé sur le bus master (plage : 200 Hz à 20000 Hz).
    \item \textbf{Curseur Résonance (Q) :} Ajuste le facteur de qualité (résonance) du filtre passe-bas master.
    \item \textbf{Curseur Reverb (macroRoom) :} Variable globale de diffusion spatiale utilisable dans vos morceaux.
    \item \textbf{Curseur Volume Général :} Contrôle le gain matériel final (\texttt{GainNode}) sans distorsion numérique.
    \item \textbf{Accès en Tiroir Compact :} En mode fenêtré étroit, le visualisateur et les contrôles master s'ouvrent en superposition d'un simple clic sur \texttt{[Effets \& Visu]}.
\end{itemize}

\subsection{6. Centre de Mises à Jour Graphique Intégré}
Strudel Studio élimine totalement les dialogues \texttt{alert()} natifs du navigateur au profit d'une interface modale sobre, réactive et ergonomique :
\begin{itemize}
    \item \textbf{Vérification en Arrière-Plan (\texttt{checking}) :} Détecte automatiquement la révision distante sur le dépôt GitHub \texttt{HugoBirard/Strudel}.
    \item \textbf{État Confirmé À Jour (\texttt{up-to-date}) :} Affiche le badge de conformité avec l'identifiant court (SHA) du commit en cours d'exécution.
    \item \textbf{Mise à Jour Détectée (\texttt{update-available}) :} Affiche le nombre de commits d'avance sur le dépôt distant ainsi que le \textbf{journal complet des modifications (changelog)} avant d'appliquer la mise à niveau via le bouton \textit{"Mettre à jour maintenant"}.
    \item \textbf{Progression en Direct (\texttt{installing}) :} Affiche une jauge d'avancement animée pendant l'exécution des opérations \texttt{git pull} et \texttt{npm install}, suivie du redémarrage automatique de l'application.
    \item \textbf{Résilience \& Diagnostic (\texttt{error}) :} En cas de déconnexion réseau, un statut explicatif est proposé sans interrompre votre session musicale.
\end{itemize}

\newpage

\section{Guide Musical Pas-à-Pas pour Débutants}
La composition musicale en live-coding s'articule traditionnellement autour de couches superposées grâce à la fonction \texttt{stack()}. Voici l'ordre recommandé pour construire un morceau cohérent :

\subsection{Étape 1 : Le Rythme (Batterie \& Percussions)}
Le rythme pose l'assise du morceau. On utilise la fonction \texttt{s()} avec une banque de sons telle que la Roland TR-808 :
\begin{verbatim}
s("bd [~ bd] sd [~ cp]")
  .bank("RolandTR808")
\end{verbatim}

\subsection{Étape 2 : La Ligne de Basse (Bassline)}
La basse ancre l'harmonie et dialogue avec la grosse caisse (\textit{kick}) :
\begin{verbatim}
s("sawtooth")
  .note("<c2 [c2 eb2] g2 bb1>")
  .lpf(800)
  .decay(0.2)
  .sustain(0.1)
\end{verbatim}

\subsection{Étape 3 : L'Harmonie (Accords \& Nappes)}
Les accords installent l'émotion et l'atmosphère générale du morceau :
\begin{verbatim}
s("gm_electric_piano_1")
  .chord("<Cm7!2 Fm7 G7>")
  .room(0.5)
  .gain(0.7)
\end{verbatim}

\subsection{Étape 4 : Le Motif Mélodique ou Arpège}
Une mélodie rapide ou un arpège donne du mouvement :
\begin{verbatim}
s("triangle")
  .note("c4 eb4 g4 bb4 c5 bb4 g4 eb4")
  .fast(2)
  .decay(0.12)
  .sustain(0.02)
  .delay(0.3)
\end{verbatim}

\subsection{Étape 5 : L'Assemblage Complet avec \texttt{stack()}}
On rassemble l'ensemble des pistes au sein d'un seul bloc exécutable :
\begin{verbatim}
stack(
  // Piste 1 : Percussions
  s("bd [~ bd] [sd,cp] [~ bd]").bank("RolandTR808"),
  s("hh*8").bank("RolandTR808").gain(0.6),

  // Piste 2 : Basse
  s("sawtooth").note("<c2 [c2 eb2] g2 bb1>").lpf(1000).decay(0.18),

  // Piste 3 : Accords
  s("gm_electric_piano_1").chord("<Cm7!2 Fm7 G7>").room(0.4).gain(0.65),

  // Piste 4 : Arpège
  s("sine").note("c4 eb4 g4 c5").fast(4).decay(0.08).gain(0.5)
)
\end{verbatim}

\newpage

\section{Catalogue Exhaustif des 128 Instruments General MIDI}
Strudel intègre l'ensemble de la norme General MIDI. Pour jouer un instrument, utilisez la syntaxe :\\
\texttt{s("gm\_nom\_instrument").note("c3 e3 g3")}

\subsection{Pianos (0--7)}
\begin{tabularx}{\linewidth}{|l|X|}
\hline
\textbf{Nom Strudel} & \textbf{Description} \\
\hline
\texttt{gm\_acoustic\_grand\_piano} & Grand piano acoustique de concert \\
\texttt{gm\_bright\_acoustic\_piano} & Piano acoustique brillant \\
\texttt{gm\_electric\_grand\_piano} & Grand piano électrique CP-70 \\
\texttt{gm\_honky\_tonk\_piano} & Piano désaccordé de saloon \\
\texttt{gm\_electric\_piano\_1} & Piano électrique Rhodes \\
\texttt{gm\_electric\_piano\_2} & Piano électrique FM type DX7 \\
\texttt{gm\_harpsichord} & Clavecin baroque \\
\texttt{gm\_clavinet} & Clavinet funky \\
\hline
\end{tabularx}

\subsection{Percussions Chromatiques (8--15)}
\begin{tabularx}{\linewidth}{|l|X|}
\hline
\textbf{Nom Strudel} & \textbf{Description} \\
\hline
\texttt{gm\_celesta} & Célesta féérique \\
\texttt{gm\_glockenspiel} & Glockenspiel métallique \\
\texttt{gm\_music\_box} & Boîte à musique mécanique \\
\texttt{gm\_vibraphone} & Vibraphone de jazz avec trémolo \\
\texttt{gm\_marimba} & Marimba en bois chaud \\
\texttt{gm\_xylophone} & Xylophone percutant \\
\texttt{gm\_tubular\_bells} & Cloches tubulaires d'orchestre \\
\texttt{gm\_dulcimer} & Tympanon traditionnel \\
\hline
\end{tabularx}

\subsection{Orgues (16--23)}
\begin{tabularx}{\linewidth}{|l|X|}
\hline
\textbf{Nom Strudel} & \textbf{Description} \\
\hline
\texttt{gm\_drawbar\_organ} & Orgue Hammond B3 \\
\texttt{gm\_percussive\_organ} & Orgue percussif gospel / jazz \\
\texttt{gm\_rock\_organ} & Orgue rock saturé \\
\texttt{gm\_church\_organ} & Grand orgue d'église \\
\texttt{gm\_reed\_organ} & Harmonium / orgue à anches \\
\texttt{gm\_accordion} & Accordéon musette \\
\texttt{gm\_harmonica} & Harmonica blues \\
\texttt{gm\_tango\_accordion} & Bandoneon de tango \\
\hline
\end{tabularx}

\subsection{Guitares (24--31)}
\begin{tabularx}{\linewidth}{|l|X|}
\hline
\textbf{Nom Strudel} & \textbf{Description} \\
\hline
\texttt{gm\_acoustic\_guitar\_nylon} & Guitare classique cordes nylon \\
\texttt{gm\_acoustic\_guitar\_steel} & Guitare folk cordes acier \\
\texttt{gm\_electric\_guitar\_jazz} & Guitare demi-caisse jazz \\
\texttt{gm\_electric\_guitar\_clean} & Guitare électrique claire (Stratocaster) \\
\texttt{gm\_electric\_guitar\_muted} & Guitare étouffée (Palm-mute) \\
\texttt{gm\_overdriven\_guitar} & Guitare rock avec overdrive \\
\texttt{gm\_distortion\_guitar} & Guitare métal avec distorsion lourde \\
\texttt{gm\_guitar\_harmonics} & Harmoniques de guitare \\
\hline
\end{tabularx}

\subsection{Basses (32--39)}
\begin{tabularx}{\linewidth}{|l|X|}
\hline
\textbf{Nom Strudel} & \textbf{Description} \\
\hline
\texttt{gm\_acoustic\_bass} & Contrebasse acoustique jazz \\
\texttt{gm\_electric\_bass\_finger} & Basse électrique jouée aux doigts \\
\texttt{gm\_electric\_bass\_pick} & Basse électrique jouée au médiator \\
\texttt{gm\_fretless\_bass} & Basse fretless fluide \\
\texttt{gm\_slap\_bass\_1} & Basse slap funk percutante \\
\texttt{gm\_slap\_bass\_2} & Basse slap avec claquement marqué \\
\texttt{gm\_synth\_bass\_1} & Basse synthé analogique classique \\
\texttt{gm\_synth\_bass\_2} & Basse synthé moderne agressive \\
\hline
\end{tabularx}

\subsection{Cordes \& Orchestre (40--55)}
\begin{tabularx}{\linewidth}{|l|X|}
\hline
\textbf{Nom Strudel} & \textbf{Description} \\
\hline
\texttt{gm\_violin} & Violon solo \\
\texttt{gm\_viola} & Alto solo \\
\texttt{gm\_cello} & Violoncelle chaleureux \\
\texttt{gm\_contrabass} & Contrebasse à l'archet \\
\texttt{gm\_tremolo\_strings} & Ensemble de cordes en trémolo \\
\texttt{gm\_pizzicato\_strings} & Cordes jouées en pizzicato \\
\texttt{gm\_orchestral\_harp} & Harpe de concert \\
\texttt{gm\_timpani} & Timbales orchestrales \\
\texttt{gm\_string\_ensemble\_1} & Ensemble de cordes symphoniques \\
\texttt{gm\_string\_ensemble\_2} & Ensemble de cordes lentes (Slow strings) \\
\texttt{gm\_synth\_strings\_1} & Nappe de cordes synthé analogique \\
\texttt{gm\_synth\_strings\_2} & Nappe de cordes synthé douce \\
\texttt{gm\_choir\_aahs} & Chœur vocal classique (Aahs) \\
\texttt{gm\_voice\_oohs} & Chœur vocal doux (Oohs) \\
\texttt{gm\_synth\_voice} & Voix synthétisée vocoder \\
\texttt{gm\_orchestra\_hit} & Coup orchestral dramatique \\
\hline
\end{tabularx}

\subsection{Cuivres (Brass) (56--63)}
\begin{tabularx}{\linewidth}{|l|X|}
\hline
\textbf{Nom Strudel} & \textbf{Description} \\
\hline
\texttt{gm\_trumpet} & Trompette éclatante \\
\texttt{gm\_trombone} & Trombone à coulisse \\
\texttt{gm\_tuba} & Tuba puissant \\
\texttt{gm\_muted\_trumpet} & Trompette avec sourdine jazz \\
\texttt{gm\_french\_horn} & Cor d'harmonie symphonique \\
\texttt{gm\_brass\_section} & Section complète de cuivres cuivrés \\
\texttt{gm\_synth\_brass\_1} & Cuivres synthétiques 80s (Jump style) \\
\texttt{gm\_synth\_brass\_2} & Cuivres synthétiques tranchants \\
\hline
\end{tabularx}

\subsection{Bois (Reed \& Pipe) (64--79)}
\begin{tabularx}{\linewidth}{|l|X|}
\hline
\textbf{Nom Strudel} & \textbf{Description} \\
\hline
\texttt{gm\_soprano\_sax} & Saxophone soprano \\
\texttt{gm\_alto\_sax} & Saxophone alto expressif \\
\texttt{gm\_tenor\_sax} & Saxophone ténor chaud \\
\texttt{gm\_baritone\_sax} & Saxophone baryton grave \\
\texttt{gm\_oboe} & Hautbois pastoral \\
\texttt{gm\_english\_horn} & Cor anglais mélancolique \\
\texttt{gm\_bassoon} & Basson orchestral \\
\texttt{gm\_clarinet} & Clarinette classique \\
\texttt{gm\_piccolo} & Petite flûte piccolo aiguë \\
\texttt{gm\_flute} & Grande flûte traversière \\
\texttt{gm\_recorder} & Flûte à bec \\
\texttt{gm\_pan\_flute} & Flûte de Pan \\
\texttt{gm\_blown\_bottle} & Son de bouteille soufflée \\
\texttt{gm\_shakuhachi} & Flûte traditionnelle japonaise \\
\texttt{gm\_whistle} & Sifflet aigu \\
\texttt{gm\_ocarina} & Ocarina en terre cuite \\
\hline
\end{tabularx}

\subsection{Synth Leads (80--87) \& Synth Pads (88--95)}
\begin{tabularx}{\linewidth}{|l|X|}
\hline
\textbf{Nom Strudel} & \textbf{Description} \\
\hline
\texttt{gm\_lead\_1\_square} & Lead onde carrée classique chiptune \\
\texttt{gm\_lead\_2\_sawtooth} & Lead onde dent de scie tranchant \\
\texttt{gm\_lead\_3\_calliope} & Lead vapeur rétro calliope \\
\texttt{gm\_lead\_4\_chiff} & Lead percutant avec souffle \\
\texttt{gm\_lead\_5\_charang} & Lead guitare synthétique moderne \\
\texttt{gm\_lead\_6\_voice} & Lead vocal chantant \\
\texttt{gm\_lead\_7\_fifths} & Lead avec quinte intégrée \\
\texttt{gm\_lead\_8\_bass\_and\_lead} & Lead puissant doublé à l'octave \\
\texttt{gm\_pad\_1\_new\_age} & Nappe new age planante \\
\texttt{gm\_pad\_2\_warm} & Nappe analogique chaude \\
\texttt{gm\_pad\_3\_polysynth} & Nappe polysynthé classique \\
\texttt{gm\_pad\_4\_choir} & Nappe de chœur céleste \\
\texttt{gm\_pad\_5\_bowed} & Nappe de cordes frottées \\
\texttt{gm\_pad\_6\_metallic} & Nappe aux reflets métalliques \\
\texttt{gm\_pad\_7\_halo} & Nappe lumineuse avec halo résonant \\
\texttt{gm\_pad\_8\_sweep} & Nappe avec balayage de filtre lent \\
\hline
\end{tabularx}

\subsection{Synth FX \& World (96--111)}
\begin{tabularx}{\linewidth}{|l|X|}
\hline
\textbf{Nom Strudel} & \textbf{Description} \\
\hline
\texttt{gm\_fx\_1\_rain} & Effet sonore d'averse \\
\texttt{gm\_fx\_2\_soundtrack} & Son d'ambiance cinématographique \\
\texttt{gm\_fx\_3\_crystal} & Cristaux carillonnants scintillants \\
\texttt{gm\_fx\_4\_atmosphere} & Nappe de science-fiction spatiale \\
\texttt{gm\_fx\_5\_brightness} & Sonorité lumineuse riche en harmoniques \\
\texttt{gm\_fx\_6\_goblins} & Effet mystérieux et fantastique \\
\texttt{gm\_fx\_7\_echoes} & Sons avec échos multiples \\
\texttt{gm\_fx\_8\_sci\_fi} & Sons de vaisseaux et technologies SF \\
\texttt{gm\_sitar} & Cithare indienne traditionnelle \\
\texttt{gm\_banjo} & Banjo américain country \\
\texttt{gm\_shamisen} & Luth japonais trois cordes \\
\texttt{gm\_koto} & Harpe japonaise à chevalets mobiles \\
\texttt{gm\_kalimba} & Piano à pouces africain \\
\texttt{gm\_bag\_pipe} & Cornemuse écossaise \\
\texttt{gm\_fiddle} & Violon traditionnel folk \\
\texttt{gm\_shanai} & Hautbois indien festif \\
\hline
\end{tabularx}

\subsection{Percussions \& Sons Spéciaux (112--127)}
\begin{tabularx}{\linewidth}{|l|X|}
\hline
\textbf{Nom Strudel} & \textbf{Description} \\
\hline
\texttt{gm\_tinkle\_bell} & Clochette tintante \\
\texttt{gm\_agogo} & Cloches agogô brésiliennes \\
\texttt{gm\_steel\_drums} & Steel drum caribéen (Pans) \\
\texttt{gm\_woodblock} & Bloc de bois sec \\
\texttt{gm\_taiko\_drum} & Gros tambour taiko japonais \\
\texttt{gm\_melodic\_tom} & Tom de batterie mélodique accordable \\
\texttt{gm\_synth\_drum} & Tambour synthétique électronique 80s \\
\texttt{gm\_reverse\_cymbal} & Cymbale inversée de transition \\
\texttt{gm\_guitar\_fret\_noise} & Bruit de frette et frottement de corde \\
\texttt{gm\_breath\_noise} & Bruit d'aspiration et de souffle humain \\
\texttt{gm\_seashore} & Vagues et ressenti de bord de mer \\
\texttt{gm\_bird\_tweet} & Gazouillis d'oiseaux dans la nature \\
\texttt{gm\_telephone\_ring} & Sonnerie de téléphone vintage \\
\texttt{gm\_helicopter} & Bruit de pales d'hélicoptère \\
\texttt{gm\_applause} & Applaudissements et acclamations de foule \\
\texttt{gm\_gunshot} & Coup de feu sec et retentissant \\
\hline
\end{tabularx}

\newpage

\section{Banques de Percussions \& Synthétiseurs}

\subsection{Banques de Boîtes à Rythmes Disponibles}
Appliquez une banque avec \texttt{.bank("NomBanque")} :
\begin{itemize}
    \item \texttt{RolandTR808} : Le son intemporel du hip-hop, de la trap et de l'électro (grosses basses 808, claps distinctifs, cowbells).
    \item \texttt{RolandTR909} : La référence de la techno et de la house music (kick percutant, charlestons ouverts ouverts et caisses claires claquantes).
    \item \texttt{RolandTR707} : Sons numériques 12-bit typiques de la synthwave et de la pop des années 80.
    \item \texttt{LinnDrum} : La boîte à rythmes mythique utilisée par Prince, Michael Jackson et New Order.
    \item \texttt{AkaiMPC60} : Percussions hip-hop old school chaleureuses avec grain vintage.
\end{itemize}

\subsection{Noms des Éléments de Batterie (Drum Hits)}
\begin{tabular}{|l|l|}
\hline
\textbf{Symbole} & \textbf{Élément} \\
\hline
\texttt{bd} & Bass Drum / Kick (Grosse caisse) \\
\texttt{sd} & Snare Drum (Caisse claire) \\
\texttt{hh} & Closed Hi-Hat (Charleston fermé) \\
\texttt{oh} & Open Hi-Hat (Charleston ouvert) \\
\texttt{cp} & Clap (Claquement de mains) \\
\texttt{cb} & Cowbell (Cloche à vache) \\
\texttt{rim} & Rimshot (Frappe sur le cercle) \\
\texttt{cr} & Crash Cymbal (Cymbale crash) \\
\texttt{rd} & Ride Cymbal (Cymbale ride) \\
\texttt{ht}, \texttt{mt}, \texttt{lt} & High, Mid, Low Toms \\
\hline
\end{tabular}

\subsection{Synthétiseurs Superdough Internes (100\% Hors-Ligne)}
Ces générateurs d'ondes fonctionnent immédiatement sans nécessiter aucun téléchargement de samples :
\begin{itemize}
    \item \texttt{s("sawtooth")} : Onde en dent de scie riche en harmoniques, idéale pour les basses techno et les leads saturés.
    \item \texttt{s("square")} : Onde carrée creuse et percutante, base des sons chiptune 8-bit et orgues d'arcade.
    \item \texttt{s("triangle")} : Onde triangulaire douce et pure, parfaite pour les mélodies cristallines et les flûtes synthétiques.
    \item \texttt{s("sine")} : Onde sinusoïdale pure sans aucune harmonique, base des sub-basses profondes.
\end{itemize}

\section{Guide des Effets Audio \& Filtres}
Tous les effets peuvent être chaînés en fin de pattern :

\begin{tabularx}{\linewidth}{|l|l|X|}
\hline
\textbf{Méthode} & \textbf{Valeurs types} & \textbf{Description} \\
\hline
\texttt{.lpf(freq)} & \texttt{100} à \texttt{18000} & Filtre Passe-Bas : atténue les fréquences supérieures \\
\texttt{.lpq(q)} & \texttt{1} à \texttt{15} & Résonance du filtre passe-bas \\
\texttt{.hpf(freq)} & \texttt{100} à \texttt{8000} & Filtre Passe-Haut : coupe les sous-basses indésirables \\
\texttt{.decay(sec)} & \texttt{0.05} à \texttt{2.0} & Vitesse d'extinction de la note après l'attaque \\
\texttt{.sustain(val)} & \texttt{0.0} à \texttt{1.0} & Niveau de volume maintenu tant que la note est active \\
\texttt{.gain(val)} & \texttt{0.0} à \texttt{1.2} & Ajustement individuel du volume de la piste \\
\texttt{.pan(val)} & \texttt{0.0} (G) à \texttt{1.0} (D) & Positionnement stéréophonique gauche / droite \\
\texttt{.room(val)} & \texttt{0.0} à \texttt{1.0} & Quantité de réverbération et d'espace acoustique \\
\texttt{.delay(val)} & \texttt{0.0} à \texttt{1.0} & Quantité d'écho et de répétition temporelle \\
\texttt{.delaytime(val)} & \texttt{0.125} à \texttt{0.5} & Temps de décalage de l'écho en fractions de cycle \\
\texttt{.shape(val)} & \texttt{0.0} à \texttt{0.9} & Saturation et distorsion harmonique chaleureuse \\
\texttt{.crush(bits)} & \texttt{4} à \texttt{12} & Réduction de bits numérique (Bitcrusher rétro) \\
\hline
\end{tabularx}

\newpage

\section{Guide de la Mini-Notation Rythmique TidalCycles}
La mini-notation permet d'écrire des rythmes extrêmement complexes et polyrythmiques en quelques caractères :

\begin{tabularx}{\linewidth}{|l|l|X|}
\hline
\textbf{Opérateur} & \textbf{Exemple} & \textbf{Signification} \\
\hline
\textbf{Espace} & \texttt{"c3 e3 g3"} & Joue les notes en succession équitable dans le cycle \\
\textbf{Crochets \texttt{[]}} & \texttt{"[bd sd]"} & Regroupe plusieurs événements dans la durée d'un seul temps \\
\textbf{Tilde \texttt{\~}} & \texttt{"[\~{} sd]"} & Silence (pause) d'une durée équivalente \\
\textbf{Multiplication \texttt{*}} & \texttt{"hh*4"} & Répète la note 4 fois plus vite dans le même intervalle \\
\textbf{Division \texttt{/}} & \texttt{"c3/2"} & Étire la note sur 2 cycles consécutifs \\
\textbf{Virgule \texttt{,}} & \texttt{"[bd,hh*4]"} & Joue les éléments simultanément (polyphonie) \\
\textbf{Chevrons \texttt{<>}} & \texttt{"<c3 f3 g3>"} & Alterne une note par cycle consécutif \\
\textbf{Point d'exclamation \texttt{!}} & \texttt{"c3!3"} & Répète la même note 3 fois de suite \\
\textbf{Point d'interrogation \texttt{?}} & \texttt{"hh?"} & Condition de probabilité (50\% de chance de se jouer) \\
\textbf{Euclidien \texttt{(k,n)}} & \texttt{"bd(3,8)"} & Répartit équitablement 3 frappes sur 8 subdivisions temporelles \\
\hline
\end{tabularx}

\section{Raccourcis Clavier Essentiels}
\begin{itemize}
    \item \texttt{Ctrl + Entrée} (ou \texttt{Cmd + Entrée}) : \textbf{Évaluer et Jouer} le code en direct.
    \item \texttt{Ctrl + .} (ou \texttt{Cmd + .}) : \textbf{Arrêter immédiatement} le son (\textit{Hush}).
    \item \texttt{Ctrl + S} (ou \texttt{Cmd + S}) : \textbf{Sauvegarder} le morceau dans un fichier local.
    \item \texttt{Ctrl + Z} / \texttt{Ctrl + Y} : \textbf{Annuler / Rétablir} vos modifications dans l'éditeur.
    \item \texttt{Tab} : Indentation automatique du code Strudel.
\end{itemize}

\vfill
\begin{center}
\small \textcopyright\ 2026 Hugo Birard — Strudel Studio Desktop. Licence Libre AGPL-3.0.
\end{center}

\end{document}
